% TOP_PROBLEM: {"id": 30006689, "title": "Prove the Penrose inequality or present a counterexample", "category_id": 16, "category": {"id": 16, "name": "physics", "display_name": "Mathematical Physics", "description": "Problems at the intersection of mathematics and physics.", "slug": "physics", "order_index": 16, "created_at": "2026-07-31T15:26:25.671Z"}, "status": "open"}
% FIRST_POSTED: 2026-09-25
% ENTRY_KIND: statement_only
% !TeX program = lualatex
% Definition and sources only; this file is not an LLM solution attempt.
\documentclass[11pt]{article}
\usepackage[margin=25mm]{geometry}
\usepackage{fontspec}
\setmainfont{Latin Modern Roman}
\usepackage{amsmath,amssymb,mathtools}
\usepackage{unicode-math}
\setmathfont{Latin Modern Math}
\usepackage{xurl,hyperref}
\hypersetup{colorlinks=true,urlcolor=blue}
\setcounter{secnumdepth}{0}
\setlength{\parindent}{0pt}
\setlength{\parskip}{5pt}
\setlength{\emergencystretch}{3em}
\begin{document}
\section*{\texorpdfstring{Prove the Penrose inequality or present a counterexample}{Prove the Penrose inequality or present a counterexample}}\label{30006689}
\subsection{Definitions and mathematical statement}
An asymptotically flat initial data set \((M^3,g,k)\) has a Riemannian metric and a symmetric second fundamental form. On its specified Euclidean end, take
\[
g-\delta=O_2(r^{-q}),\qquad k=O_1(r^{-1-q}),\qquad q>\tfrac12,
\]
with finite ADM limits. In physical units \(G=c=1\), the constraint densities satisfy
\[
16\pi\mu=R_g+(\operatorname{tr}_gk)^2-|k|_g^2,
\qquad8\pi J=\operatorname{div}_g(k-(\operatorname{tr}_gk)g),
\qquad\mu\ge|J|_g.
\]
For a two-sided surface \(\Sigma\), choose the normal toward the designated end. A future marginally outer trapped surface has
\[
\theta_+(\Sigma)=H_\Sigma+\operatorname{tr}_\Sigma k=0.
\]
The cited formulation [S1] compares the ADM energy
\[
E_{\rm ADM}(g)=\frac{1}{16\pi}\lim_{R\to\infty}\int_{S_R}
(\partial_jg_{ij}-\partial_i g_{jj})n_\delta^i\,dA_\delta
\]
with the actual area \(A_\Sigma=\operatorname{Area}_g(\Sigma)\) of the specified horizon or trapped surface:
\[
E_{\rm ADM}(g)\ge\sqrt{\frac{A_\Sigma}{16\pi}}.
\]
Its stated rigidity identifies equality with original data induced by a Schwarzschild slice and a unique connected outermost horizon. The full spacetime question allows \(k\ne0\). Source [S1, Theorem B] asserts a conditional result for an outermost stable MOTS with \(\operatorname{tr}_\Sigma k\ge0\), or its stated distributional KKT alternative; its general trapped-surface extension adds further conditions. These are hypotheses of that author's claim, rather than an independently verified resolution.

There is a distinct minimum-enclosing-area formulation. If \(\Omega\) is the chosen exterior and \(\Gamma\) ranges over its full enclosing cuts, put
\[
A_{\min}=\inf_\Gamma\operatorname{Area}_g(\Gamma),\qquad
m_{\rm inv}=\sqrt{E_{\rm ADM}^2-|P_{\rm ADM}|^2}.
\]
The corresponding inequality is
\[
m_{\rm inv}\ge\sqrt{\frac{A_{\min}}{16\pi}}.
\]
Here every cut component and every portion coinciding with the inner boundary contributes to area. The cut infimum can be smaller than \(A_\Sigma\); these formulations therefore require separate hypotheses and scope checks. In an outer area-minimizing horizon class, \(A_{\min}=A_\Sigma\).

\subsection{Short English statement}
Does the mass of an asymptotically flat spacetime initial state always dominate the area-based lower bound from its outermost horizon?
\subsection{Sources}
\begin{itemize}
\item[S1]D. Xu, 'The Spacetime Penrose Inequality: Conditional Results for Stable MOTS and General Trapped Surfaces,' arXiv:2512.04137v5.
\url{https://arxiv.org/abs/2512.04137v5}.
\textit{Formulation source.}
\item[S2]The spacetime Penrose inequality under a quasi final state hypothesis.
\url{https://arxiv.org/abs/2605.18730}.
\textit{Further reference; not a proof endorsement.}
\item[S3]Proof of the Spacetime Penrose Inequality With Suboptimal Constant in the Asymptotically Flat and Asymptotically Hyperboloidal Regimes.
\url{https://arxiv.org/abs/2504.10641}.
\textit{Further reference; not a proof endorsement.}
\end{itemize}
\end{document}
